\documentclass[11pt]{article}
\usepackage[margin=1in]{geometry}
\usepackage{booktabs}
\usepackage{caption}
\usepackage{amsmath}
\usepackage{float}
\usepackage{graphicx}
\usepackage{adjustbox}

\begin{document}

\section*{Asymmetric Demand Multipliers by Shock Size: v4}

This note reports the v4 pooled positive/negative specification. Like v3, v4
keeps positive and negative shocks in the same monthly cross-section and
estimates six sign-specific bin slopes in one regression. The difference is
that v4 rebuilds the positive and negative bins directly from raw demand using
sign-specific standard deviations within each demand type and month.

For each demand type and month, define
\[
\sigma^+_t = \operatorname{sd}(d_{n,t} \mid d_{n,t} > 0), \qquad
\sigma^-_t = \operatorname{sd}(d_{n,t} \mid d_{n,t} < 0).
\]
The six regressors are
\begin{align*}
d^{(1),+}_{n,t} &= d_{n,t}\mathbf{1}\{0 < d_{n,t} < \sigma^+_t\}, \\
d^{(2),+}_{n,t} &= d_{n,t}\mathbf{1}\{\sigma^+_t \le d_{n,t} < 2\sigma^+_t\}, \\
d^{(3),+}_{n,t} &= d_{n,t}\mathbf{1}\{d_{n,t} \ge 2\sigma^+_t\},
\end{align*}
and
\begin{align*}
d^{(1),-}_{n,t} &= d_{n,t}\mathbf{1}\{-\sigma^-_t < d_{n,t} < 0\}, \\
d^{(2),-}_{n,t} &= d_{n,t}\mathbf{1}\{-2\sigma^-_t < d_{n,t} \le -\sigma^-_t\}, \\
d^{(3),-}_{n,t} &= d_{n,t}\mathbf{1}\{d_{n,t} \le -2\sigma^-_t\}.
\end{align*}
\subsection*{Static specification}

The static monthly Fama--MacBeth regression is
\[
r_{n,t} = \alpha_t
        + \sum_{k=1}^3 M_k^+ d^{(k),+}_{n,t}
        + \sum_{k=1}^3 M_k^- d^{(k),-}_{n,t}
        + \Gamma X_{n,t} + \varepsilon_{n,t}.
\]

The regression is estimated separately for BMI, FIT, and OFI demand. Columns
(1)--(3) use BMI, columns (4)--(6) use FIT, and columns (7)--(9) use OFI.
Within each demand type, the three columns correspond to no controls,
predictor controls, and predictor plus liquidity controls. BMI specifications
also include the BMI-specific control set used in the main price-impact tables.

\paragraph{Implementation detail.}
The script computes $\sigma^+_t$ and $\sigma^-_t$ within demand type and month.
This is the all-type R equivalent of the Stata construction
\texttt{bysort yyyymm} after first filtering to one demand type. No
type-month has missing positive- or negative-side standard deviations.

If a monthly cross-section omits a coefficient because the corresponding bin
has no observations, v4 keeps that monthly coefficient missing when computing
Fama--MacBeth averages. It does not replace omitted coefficients with zero.
The difference rows are Fama--MacBeth averages of monthly paired differences,
not differences of the printed coefficient averages; these can diverge when
the two coefficient rows use different nonmissing month sets.

\paragraph{Takeaway.}
The positive-demand domain continues to show strong concavity for FIT and OFI:
the multiplier is largest for small positive shocks and smallest for large
positive shocks. This is not consistent with a short-sales-constraint
interpretation under which multipliers should rise with shock size. On the
negative-demand side, OFI also shows strong concavity. FIT is flatter on the
negative side, with the largest-bin multiplier below the smallest-bin
multiplier but small and imprecise adjacent-bin differences. BMI remains noisy
because it has only 21 monthly cross-sections. Overall, v4 does not show
evidence that $M$ increases with shock size, so the results are not consistent
with leverage constraints generating larger multipliers for larger shocks.

\begin{table}[H]
\centering
\caption{v4 rebuilt shock-size bin summary statistics}
\label{tab:v4_bin_summary}
\begin{tabular}{llrrrr}
\toprule
Shock & Type & Bin & Obs & Mean $d_{n,t}$ & SD $d_{n,t}$ \\
\midrule
Negative & BMI & 1 & 3,671 & -0.001614 & 0.001317 \\
Negative & BMI & 2 & 477 & -0.007543 & 0.002594 \\
Negative & BMI & 3 & 533 & -0.016797 & 0.006574 \\
Negative & FIT & 1 & 155,449 & -0.001062 & 0.001012 \\
Negative & FIT & 2 & 49,096 & -0.004405 & 0.001770 \\
Negative & FIT & 3 & 30,160 & -0.009209 & 0.004401 \\
Negative & OFI & 1 & 235,014 & -0.012466 & 0.010089 \\
Negative & OFI & 2 & 46,995 & -0.049190 & 0.016080 \\
Negative & OFI & 3 & 24,862 & -0.121417 & 0.067350 \\
\addlinespace
Positive & BMI & 1 & 3,765 & 0.001889 & 0.001534 \\
Positive & BMI & 2 & 670 & 0.008341 & 0.002291 \\
Positive & BMI & 3 & 584 & 0.018393 & 0.006673 \\
Positive & FIT & 1 & 223,728 & 0.001369 & 0.001470 \\
Positive & FIT & 2 & 52,235 & 0.006385 & 0.002838 \\
Positive & FIT & 3 & 31,910 & 0.014020 & 0.008255 \\
Positive & OFI & 1 & 163,014 & 0.010444 & 0.009304 \\
Positive & OFI & 2 & 37,776 & 0.038842 & 0.018288 \\
Positive & OFI & 3 & 20,083 & 0.094069 & 0.056545 \\
\bottomrule
\end{tabular}
\end{table}

\begin{table}[H]
\centering
\caption{v4 rebuilt-bin pooled asymmetric multipliers by shock-size bin}
\label{tab:v4_asym_multipliers}
\begin{adjustbox}{max width=\textwidth}
\begingroup
\renewcommand{\vspace}[1]{}
\begin{tabular}{lccccccccc}
  \hline \multicolumn{10}{c}{Panel A: Positive shocks} \\
  \hline
  & \multicolumn{3}{c}{BMI} & \multicolumn{3}{c}{FIT} & \multicolumn{3}{c}{OFI} \\
  \cmidrule(l){2-4} \cmidrule(l){5-7} \cmidrule(l){8-10}
  & (1) & (2) & (3) & (4) & (5) & (6) & (7) & (8) & (9) \\
  $M_{\{0 < d_{n,t} < \sigma^+\}}$ & $0.39$ & $1.09$ & $1.46$ & $10.48^{***}$ & $8.25^{***}$ & $7.32^{***}$ & $3.79^{***}$ & $3.87^{***}$ & $4.34^{***}$ \\
   & $(1.61)$ & $(1.45)$ & $(1.12)$ & $(1.69)$ & $(1.35)$ & $(1.06)$ & $(0.25)$ & $(0.21)$ & $(0.20)$ \\
  $M_{\{d_{n,t} \in [\sigma^+, 2\sigma^+)\}}$ & $0.29$ & $0.33$ & $0.62$ & $7.90^{***}$ & $6.52^{***}$ & $6.09^{***}$ & $2.76^{***}$ & $2.79^{***}$ & $3.06^{***}$ \\
   & $(0.66)$ & $(0.79)$ & $(0.75)$ & $(1.06)$ & $(0.90)$ & $(0.81)$ & $(0.21)$ & $(0.18)$ & $(0.15)$ \\
  $M_{\{d_{n,t} \ge 2\sigma^+\}}$ & $0.78^{**}$ & $0.63^{*}$ & $0.69^{*}$ & $4.75^{***}$ & $4.21^{***}$ & $4.14^{***}$ & $1.78^{***}$ & $1.78^{***}$ & $1.92^{***}$ \\
   & $(0.38)$ & $(0.36)$ & $(0.39)$ & $(0.57)$ & $(0.52)$ & $(0.50)$ & $(0.17)$ & $(0.15)$ & $(0.13)$ \\
  \vspace{5pt}
  Predictor controls & N & Y & Y & N & Y & Y & N & Y & Y \\
  Liquidity controls & N & N & Y & N & N & Y & N & N & Y \\
  \vspace{5pt}
  Obs & 9,908 & 9,908 & 9,908 & 542,578 & 542,578 & 542,578 & 527,744 & 527,744 & 527,744 \\
  $R^2$ & 0.060 & 0.149 & 0.175 & 0.013 & 0.067 & 0.082 & 0.062 & 0.113 & 0.134 \\
  Marginal $R^2(d_{n,t})$ & 0.028 & 0.021 & 0.019 & 0.013 & 0.008 & 0.007 & 0.062 & 0.053 & 0.055 \\
  \hline
  $M_{\{d_{n,t} \in [\sigma^+, 2\sigma^+)\}} - M_{\{0 < d_{n,t} < \sigma^+\}}$ & $-0.10$ & $-0.76$ & $-0.84$ & $-2.57^{**}$ & $-1.73^{*}$ & $-1.22$ & $-1.04^{***}$ & $-1.08^{***}$ & $-1.28^{***}$ \\
   & $(1.56)$ & $(1.44)$ & $(1.16)$ & $(1.20)$ & $(1.01)$ & $(0.88)$ & $(0.12)$ & $(0.11)$ & $(0.11)$ \\
  $M_{\{d_{n,t} \ge 2\sigma^+\}} - M_{\{d_{n,t} \in [\sigma^+, 2\sigma^+)\}}$ & $0.49$ & $0.30$ & $0.07$ & $-3.16^{***}$ & $-2.30^{***}$ & $-1.96^{**}$ & $-0.97^{***}$ & $-1.01^{***}$ & $-1.14^{***}$ \\
   & $(0.57)$ & $(0.66)$ & $(0.61)$ & $(0.89)$ & $(0.80)$ & $(0.76)$ & $(0.12)$ & $(0.11)$ & $(0.11)$ \\
  $M_{\{d_{n,t} \ge 2\sigma^+\}} - M_{\{0 < d_{n,t} < \sigma^+\}}$ & $0.39$ & $-0.47$ & $-0.77$ & $-5.73^{***}$ & $-4.03^{***}$ & $-3.18^{***}$ & $-2.01^{***}$ & $-2.09^{***}$ & $-2.42^{***}$ \\
   & $(1.62)$ & $(1.47)$ & $(1.11)$ & $(1.49)$ & $(1.23)$ & $(1.01)$ & $(0.16)$ & $(0.14)$ & $(0.14)$ \\
  \vspace{5pt}
  \hline \multicolumn{10}{c}{Panel B: Negative shocks} \\
  \hline
  & \multicolumn{3}{c}{BMI} & \multicolumn{3}{c}{FIT} & \multicolumn{3}{c}{OFI} \\
  \cmidrule(l){2-4} \cmidrule(l){5-7} \cmidrule(l){8-10}
  & (1) & (2) & (3) & (4) & (5) & (6) & (7) & (8) & (9) \\
  $M_{\{-\sigma^- < d_{n,t} < 0\}}$ & $6.18^{***}$ & $3.59^{**}$ & $3.14^{**}$ & $5.90^{**}$ & $5.04^{**}$ & $3.60^{***}$ & $2.69^{***}$ & $2.73^{***}$ & $2.93^{***}$ \\
   & $(1.99)$ & $(1.60)$ & $(1.44)$ & $(2.68)$ & $(1.98)$ & $(1.40)$ & $(0.16)$ & $(0.13)$ & $(0.14)$ \\
  $M_{\{d_{n,t} \in (-2\sigma^-, -\sigma^-]\}}$ & $2.55^{**}$ & $1.46^{*}$ & $1.28^{*}$ & $4.57^{**}$ & $3.81^{***}$ & $3.32^{***}$ & $1.79^{***}$ & $1.78^{***}$ & $1.85^{***}$ \\
   & $(1.03)$ & $(0.78)$ & $(0.77)$ & $(1.85)$ & $(1.46)$ & $(1.16)$ & $(0.13)$ & $(0.10)$ & $(0.09)$ \\
  $M_{\{d_{n,t} \le -2\sigma^-\}}$ & $1.39^{***}$ & $1.11^{***}$ & $1.08^{***}$ & $3.94^{***}$ & $3.46^{***}$ & $3.08^{***}$ & $0.66^{***}$ & $0.67^{***}$ & $0.70^{***}$ \\
   & $(0.44)$ & $(0.37)$ & $(0.32)$ & $(0.91)$ & $(0.66)$ & $(0.54)$ & $(0.07)$ & $(0.06)$ & $(0.06)$ \\
  \vspace{5pt}
  Predictor controls & N & Y & Y & N & Y & Y & N & Y & Y \\
  Liquidity controls & N & N & Y & N & N & Y & N & N & Y \\
  \vspace{5pt}
  Obs & 9,908 & 9,908 & 9,908 & 542,578 & 542,578 & 542,578 & 527,744 & 527,744 & 527,744 \\
  $R^2$ & 0.060 & 0.149 & 0.175 & 0.013 & 0.067 & 0.082 & 0.062 & 0.113 & 0.134 \\
  Marginal $R^2(d_{n,t})$ & 0.028 & 0.021 & 0.019 & 0.013 & 0.008 & 0.007 & 0.062 & 0.053 & 0.055 \\
  \hline
  $M_{\{d_{n,t} \in (-2\sigma^-, -\sigma^-]\}} - M_{\{-\sigma^- < d_{n,t} < 0\}}$ & $-3.63^{***}$ & $-2.13$ & $-1.85$ & $-1.33$ & $-1.23$ & $-0.28$ & $-0.90^{***}$ & $-0.95^{***}$ & $-1.07^{***}$ \\
   & $(1.39)$ & $(1.39)$ & $(1.30)$ & $(1.74)$ & $(1.33)$ & $(1.19)$ & $(0.08)$ & $(0.08)$ & $(0.09)$ \\
  $M_{\{d_{n,t} \le -2\sigma^-\}} - M_{\{d_{n,t} \in (-2\sigma^-, -\sigma^-]\}}$ & $-1.16$ & $-0.35$ & $-0.21$ & $-0.63$ & $-0.35$ & $-0.24$ & $-1.13^{***}$ & $-1.11^{***}$ & $-1.15^{***}$ \\
   & $(0.97)$ & $(0.74)$ & $(0.73)$ & $(1.26)$ & $(1.15)$ & $(0.99)$ & $(0.09)$ & $(0.08)$ & $(0.07)$ \\
  $M_{\{d_{n,t} \le -2\sigma^-\}} - M_{\{-\sigma^- < d_{n,t} < 0\}}$ & $-4.79^{**}$ & $-2.48$ & $-2.06$ & $-1.96$ & $-1.58$ & $-0.52$ & $-2.03^{***}$ & $-2.05^{***}$ & $-2.22^{***}$ \\
   & $(1.98)$ & $(1.60)$ & $(1.43)$ & $(2.11)$ & $(1.58)$ & $(1.13)$ & $(0.12)$ & $(0.11)$ & $(0.12)$ \\
  \vspace{5pt}
  \hline
\end{tabular}

\endgroup
\end{adjustbox}
\end{table}

\paragraph{Controlled-specification highlights.}
In the most controlled specification, positive FIT multipliers fall from 7.32
in the smallest bin to 4.14 in the largest bin; positive OFI multipliers fall
from 4.34 to 1.92. Negative OFI multipliers fall from 2.93 to 0.70, with a
largest-minus-smallest difference of -2.22. Negative FIT falls from 3.60 to
3.08 between the smallest and largest bins, but the differences are imprecise.
BMI coefficients are not the main evidence because the sample has only 21
months and the estimates are noisy.

\begin{figure}[H]
\centering
\includegraphics[width=0.32\textwidth]{plots/v4/asym_v4_BMI.png}
\includegraphics[width=0.32\textwidth]{plots/v4/asym_v4_FIT.png}
\includegraphics[width=0.32\textwidth]{plots/v4/asym_v4_OFI.png}
\caption{v4 static pairwise multiplier differences, most controlled specification}
\label{fig:v4_static_tight_controls}
\end{figure}

\subsection*{Dynamic specification}

The dynamic v4 specification applies the same pooled six-variable idea to the
dynamic regression input. For each horizon $h$, sign and bin membership are
based on the lagged cumulative demand shock, denoted $c^{(h)}_{n,t}$ in this
note. The sign-specific standard deviations are computed within demand type,
horizon, and month:
\[
\sigma^{+}_{t,h} = \operatorname{sd}(c^{(h)}_{n,t} \mid c^{(h)}_{n,t} > 0),
\qquad
\sigma^{-}_{t,h} = \operatorname{sd}(c^{(h)}_{n,t} \mid c^{(h)}_{n,t} < 0).
\]
Bin membership uses these lagged cumulative shocks, while the regressor value
inside each bin remains the contemporaneous demand shock $d_{n,t}$, matching
the existing non-asymmetric dynamic regression architecture:
\[
r_{n,t} = \alpha_t
        + \sum_{k=1}^3 M_{h,k}^+ d_{n,t}\mathbf{1}\{c^{(h)}_{n,t}
          \in B^{+,k}_{t,h}\}
        + \sum_{k=1}^3 M_{h,k}^- d_{n,t}\mathbf{1}\{c^{(h)}_{n,t}
          \in B^{-,k}_{t,h}\}
        + \Gamma_h X_{n,t} + \varepsilon_{n,t}.
\]
The table below reports the 4-lag horizon for FIT and OFI. The dynamic v4
balanced-month filter drops zero type-horizon-month cells, so the balanced and
unbalanced dynamic v4 results are identical.

\begin{table}[H]
\centering
\caption{v4 dynamic rebuilt 4-lag bin summary statistics}
\label{tab:v4_dynamic_bin_summary_4lag}
\begin{adjustbox}{max width=\textwidth}
\begin{tabular}{lllrrrrr}
\toprule
Shock & Type & Bin & Obs & Mean $c^{(4)}_{n,t}$ & SD $c^{(4)}_{n,t}$ & Mean $d_{n,t}$ & SD $d_{n,t}$ \\
\midrule
Negative & FIT & 1 & 114,326 & -0.003445 & 0.003020 & -0.000283 & 0.004099 \\
Negative & FIT & 2 & 40,622 & -0.013349 & 0.004937 & -0.001602 & 0.005754 \\
Negative & FIT & 3 & 25,130 & -0.027758 & 0.011936 & -0.003641 & 0.007187 \\
Negative & OFI & 1 & 185,111 & -0.032991 & 0.026626 & -0.007569 & 0.027950 \\
Negative & OFI & 2 & 45,527 & -0.118758 & 0.046659 & -0.021215 & 0.042954 \\
Negative & OFI & 3 & 23,831 & -0.290753 & 0.185539 & -0.047018 & 0.071998 \\
\addlinespace
Positive & FIT & 1 & 193,499 & 0.004533 & 0.004299 & 0.000972 & 0.003832 \\
Positive & FIT & 2 & 54,975 & 0.018380 & 0.007660 & 0.002933 & 0.006293 \\
Positive & FIT & 3 & 32,623 & 0.040165 & 0.021259 & 0.005684 & 0.009530 \\
Positive & OFI & 1 & 121,294 & 0.032871 & 0.033268 & 0.003089 & 0.028916 \\
Positive & OFI & 2 & 31,142 & 0.116028 & 0.076267 & 0.017751 & 0.043433 \\
Positive & OFI & 3 & 16,791 & 0.281708 & 0.207627 & 0.046401 & 0.074548 \\
\bottomrule
\end{tabular}

\end{adjustbox}
\end{table}

\begin{table}[H]
\centering
\caption{v4 dynamic rebuilt-bin pooled asymmetric multipliers, 4-lag horizon}
\label{tab:v4_dynamic_asym_multipliers_4lag}
\begin{adjustbox}{max width=\textwidth}
\begingroup
\renewcommand{\vspace}[1]{}
\begin{tabular}{lcccccc}
  \hline \multicolumn{7}{c}{Panel A: Positive lagged cumulative shocks} \\
  \hline
  & \multicolumn{3}{c}{FIT} & \multicolumn{3}{c}{OFI} \\
  \cmidrule(l){2-4} \cmidrule(l){5-7}
  & (1) & (2) & (3) & (4) & (5) & (6) \\
  $M_{\{0 < c^{(4)}_{n,t} < \sigma^+\}}$ & $4.89^{***}$ & $4.09^{***}$ & $3.97^{***}$ & $1.55^{***}$ & $1.54^{***}$ & $1.63^{***}$ \\
   & $(0.52)$ & $(0.37)$ & $(0.34)$ & $(0.10)$ & $(0.10)$ & $(0.10)$ \\
  $M_{\{c^{(4)}_{n,t} \in [\sigma^+, 2\sigma^+)\}}$ & $3.83^{***}$ & $3.35^{***}$ & $3.33^{***}$ & $0.98^{***}$ & $0.98^{***}$ & $1.09^{***}$ \\
   & $(0.55)$ & $(0.43)$ & $(0.41)$ & $(0.12)$ & $(0.11)$ & $(0.11)$ \\
  $M_{\{c^{(4)}_{n,t} \ge 2\sigma^+\}}$ & $2.42^{***}$ & $2.34^{***}$ & $2.29^{***}$ & $0.46^{***}$ & $0.47^{***}$ & $0.57^{***}$ \\
   & $(0.51)$ & $(0.42)$ & $(0.42)$ & $(0.12)$ & $(0.11)$ & $(0.11)$ \\
  \vspace{5pt}
  Predictor controls & N & Y & Y & N & Y & Y \\
  Liquidity controls & N & N & Y & N & N & Y \\
  \vspace{5pt}
  Obs & 461,175 & 461,175 & 461,175 & 423,696 & 423,696 & 423,696 \\
  $R^2$ & 0.012 & 0.073 & 0.087 & 0.054 & 0.112 & 0.130 \\
  Marginal $R^2(d_{n,t})$ & 0.012 & 0.007 & 0.006 & 0.054 & 0.045 & 0.045 \\
  \hline
  $M_{\{c^{(4)}_{n,t} \in [\sigma^+, 2\sigma^+)\}} - M_{\{0 < c^{(4)}_{n,t} < \sigma^+\}}$ & $-1.06^{**}$ & $-0.74^{*}$ & $-0.64^{*}$ & $-0.57^{***}$ & $-0.56^{***}$ & $-0.54^{***}$ \\
   & $(0.46)$ & $(0.40)$ & $(0.39)$ & $(0.09)$ & $(0.09)$ & $(0.09)$ \\
  $M_{\{c^{(4)}_{n,t} \ge 2\sigma^+\}} - M_{\{c^{(4)}_{n,t} \in [\sigma^+, 2\sigma^+)\}}$ & $-1.41^{***}$ & $-1.00^{**}$ & $-1.04^{**}$ & $-0.52^{***}$ & $-0.51^{***}$ & $-0.52^{***}$ \\
   & $(0.52)$ & $(0.42)$ & $(0.42)$ & $(0.11)$ & $(0.11)$ & $(0.10)$ \\
  $M_{\{c^{(4)}_{n,t} \ge 2\sigma^+\}} - M_{\{0 < c^{(4)}_{n,t} < \sigma^+\}}$ & $-2.47^{***}$ & $-1.75^{***}$ & $-1.68^{***}$ & $-1.10^{***}$ & $-1.07^{***}$ & $-1.06^{***}$ \\
   & $(0.53)$ & $(0.40)$ & $(0.39)$ & $(0.13)$ & $(0.12)$ & $(0.11)$ \\
  \vspace{5pt}
  \hline \multicolumn{7}{c}{Panel B: Negative lagged cumulative shocks} \\
  \hline
  & \multicolumn{3}{c}{FIT} & \multicolumn{3}{c}{OFI} \\
  \cmidrule(l){2-4} \cmidrule(l){5-7}
  & (1) & (2) & (3) & (4) & (5) & (6) \\
  $M_{\{-\sigma^- < c^{(4)}_{n,t} < 0\}}$ & $4.99^{***}$ & $4.10^{***}$ & $3.97^{***}$ & $1.96^{***}$ & $1.91^{***}$ & $1.95^{***}$ \\
   & $(0.58)$ & $(0.46)$ & $(0.42)$ & $(0.09)$ & $(0.09)$ & $(0.09)$ \\
  $M_{\{c^{(4)}_{n,t} \in (-2\sigma^-, -\sigma^-]\}}$ & $4.27^{***}$ & $3.52^{***}$ & $3.42^{***}$ & $1.76^{***}$ & $1.71^{***}$ & $1.79^{***}$ \\
   & $(0.57)$ & $(0.47)$ & $(0.44)$ & $(0.10)$ & $(0.09)$ & $(0.10)$ \\
  $M_{\{c^{(4)}_{n,t} \le -2\sigma^-\}}$ & $2.21^{***}$ & $1.64^{***}$ & $1.54^{***}$ & $1.27^{***}$ & $1.25^{***}$ & $1.35^{***}$ \\
   & $(0.53)$ & $(0.43)$ & $(0.41)$ & $(0.09)$ & $(0.08)$ & $(0.08)$ \\
  \vspace{5pt}
  Predictor controls & N & Y & Y & N & Y & Y \\
  Liquidity controls & N & N & Y & N & N & Y \\
  \vspace{5pt}
  Obs & 461,175 & 461,175 & 461,175 & 423,696 & 423,696 & 423,696 \\
  $R^2$ & 0.012 & 0.073 & 0.087 & 0.054 & 0.112 & 0.130 \\
  Marginal $R^2(d_{n,t})$ & 0.012 & 0.007 & 0.006 & 0.054 & 0.045 & 0.045 \\
  \hline
  $M_{\{c^{(4)}_{n,t} \in (-2\sigma^-, -\sigma^-]\}} - M_{\{-\sigma^- < c^{(4)}_{n,t} < 0\}}$ & $-0.72$ & $-0.58$ & $-0.55$ & $-0.21^{**}$ & $-0.19^{**}$ & $-0.16^{**}$ \\
   & $(0.55)$ & $(0.49)$ & $(0.48)$ & $(0.08)$ & $(0.08)$ & $(0.07)$ \\
  $M_{\{c^{(4)}_{n,t} \le -2\sigma^-\}} - M_{\{c^{(4)}_{n,t} \in (-2\sigma^-, -\sigma^-]\}}$ & $-2.07^{***}$ & $-1.88^{***}$ & $-1.88^{***}$ & $-0.49^{***}$ & $-0.46^{***}$ & $-0.45^{***}$ \\
   & $(0.54)$ & $(0.49)$ & $(0.49)$ & $(0.08)$ & $(0.08)$ & $(0.08)$ \\
  $M_{\{c^{(4)}_{n,t} \le -2\sigma^-\}} - M_{\{-\sigma^- < c^{(4)}_{n,t} < 0\}}$ & $-2.79^{***}$ & $-2.47^{***}$ & $-2.43^{***}$ & $-0.70^{***}$ & $-0.65^{***}$ & $-0.60^{***}$ \\
   & $(0.61)$ & $(0.55)$ & $(0.53)$ & $(0.10)$ & $(0.09)$ & $(0.08)$ \\
  \vspace{5pt}
  \hline
\end{tabular}

\endgroup
\end{adjustbox}
\end{table}

\paragraph{Dynamic controlled-specification highlights.}
In the most controlled 4-lag specification, positive FIT multipliers fall from
3.97 in the smallest lagged-cumulative-shock bin to 2.29 in the largest bin,
with a largest-minus-smallest difference of -1.68. Positive OFI multipliers
fall from 1.63 to 0.57, with a largest-minus-smallest difference of -1.06.
On the negative side, FIT falls from 3.97 to 1.54, and OFI falls from 1.95 to
1.35. The dynamic rebuilt-bin evidence therefore has the same high-level
shape as the static v4 evidence: multipliers do not rise with shock size.

\begin{figure}[H]
\centering
\includegraphics[width=0.49\textwidth]{plots_dynamic/v4_unbalanced/dynamic_asym_v4_FIT_4lag.png}
\includegraphics[width=0.49\textwidth]{plots_dynamic/v4_unbalanced/dynamic_asym_v4_OFI_4lag.png}
\caption{v4 dynamic 4-lag pairwise multiplier differences, most controlled specification}
\label{fig:v4_dynamic_tight_controls_4lag}
\end{figure}

\end{document}
